\chapter{Analysis by model}

\eyebrow{fat segmentation over the batch}

\vspace{4pt}
The \textbf{Analysis} view of a batch runs the trained models over its images and
presents the results in a workspace. This workspace holds a gallery, a sortable
table with a distribution, conformation, physical correlation, and side-by-side
comparison.

\section{Running the analysis}

A batch is opened and \textbf{Analysis} is selected. The \tele{Analyze batch}
control processes all not-yet-analysed images, with a live progress bar. Results
are saved to the batch, so re-opening does not re-run them; \tele{Reanalyze}
forces a fresh run. The \tele{CSV} control exports every result.

Each image passes through the same pipeline:

\begin{enumerate}[leftmargin=1.4em,itemsep=3pt]
  \item \textbf{Background removal.} The carcass is cut out and composited on
        black, the same conditions the fat model was trained on.
  \item \textbf{Fat segmentation.} The validated model produces a fat mask; the
        fat fraction is measured \emph{relative to the carcass surface}, not the
        whole photo.
  \item \textbf{Conformation.} The convexity measure runs on the clean silhouette
        (Station~8).
  \item \textbf{Grade (optional).} With \tele{include grade (experimental)}
        ticked, the experimental finishing grade is also computed.
\end{enumerate}

\begin{caution}[title=Why background removal matters]
The fat model was trained on cut-out carcasses on a black background. Without
removal it counts bright tiles and the operator's coat as fat (giving falsely
high values). Removal is automatic; a \tele{background removed} pill confirms it,
and the fat percentage becomes the fraction of the \emph{carcass surface}.
\end{caution}

\section{Gallery}

The gallery (Figure~\ref{fig:analysis-gallery}) shows every carcass with its fat
overlay in cyan and the fat percentage beneath, providing a visual overview of
the whole batch. Above it, batch statistics (mean, std dev, median, min, max) and
a row of status pills: \pillok{validated segmentation $\cdot$ IoU 0.92},
\pillalert{experimental grade $\cdot$ n=22}, and \tele{analysis saved to batch}.

\shot{analysis-gallery.jpeg}{Analysis $\rightarrow$ Gallery. Fat overlay (cyan) per
carcass with the fat fraction; batch statistics and the colour-coded status pills
above; the Inspector reports the batch's fat mean, std dev and range.}{analysis-gallery}

\section{Table and distribution}

The \textbf{Table + distribution} tab (Figure~\ref{fig:analysis-table}) lists the
carcasses in a sortable table (by tag or fat \%), alongside the grade and the
physical-reference columns, beside a histogram of the batch's fat distribution.

\shot{analysis-table.jpeg}{Analysis $\rightarrow$ Table + distribution. A sortable
table of per-carcass fat, grade and physical measurements, beside the batch fat
histogram.}{analysis-table}

\section{Compare}

The \textbf{Compare} tab (Figure~\ref{fig:analysis-compare}) allows up to four
carcasses to be selected and their overlays and numbers viewed side by side,
supporting calibration of perception against the model.

\shot{analysis-compare.jpeg}{Analysis $\rightarrow$ Compare. Up to four carcasses
side by side, each with its fat overlay, fat percentage and experimental grade.}{analysis-compare}

\section{Physical correlation}

The \textbf{Physical correlation} tab plots the model's fat \% against the entered
physical reference (fat thickness, loin eye area), with the Pearson $r$. This tab
supports validation once the paired dataset is large enough. It tests whether the
surface measurement predicts the physical reference measured on the animal.
